%! TEX root = part2.tex
% the main TeX file which is intended to compile, :VimtexReload after adjustment
% :h vimtex-tex-root
\documentclass[../skb.tex]{subfiles}
% \includeonlyframes{current}
% \setbeameroption{hide notes} % Only slides
% \setbeameroption{show only notes} % Only notes
% \setbeameroption{show notes on second screen=right} % Both

\begin{document}

\begin{frame}[c,mycolor=digiPH_white,mytitle=standard,light]
	\frametitle{Tugas Kelas}

	Buat dan Jelaskan \textbf{konsep bangunan tinggi} yang ingin dirancang

	Ket:
	\begin{itemize}
        \item    Kertas A3 Manual Warna
		\item    Struktur Apa
		\item    Alasan pemilihan struktur
    \end{itemize}


\end{frame}

\begin{frame}[c,mycolor=digiPH_leaf,mytitle=standard,dark]
	\includegraphics<1>[width=.9\textwidth]{../figures/cth_konsep}

	\includegraphics<2>[width=.9\textwidth]{../figures/cth_konsep_1}

	\includegraphics<3>[width=.9\textwidth]{../figures/cth_konsep_2}

\end{frame}


\begin{frame}[c,mycolor=digiPH_leaf,mytitle=standard,dark]\frametitle{Tugas Rumah}

	\begin{center}
		Buatlah gambar kerja \textbf{denah} konsep bangunan tinggi yang direncanakan.
	\end{center}

	Ket:

	\tabitem Gunakan Kop (etiket gambar) yang seragam
\end{frame}












\end{document}
